\section{Adding weighted Monge--Amp\`ere measures}

Weighted measures of two solutions in the same ordinary type can be added while retaining that type. Truncation gives bounded equations whose solutions have a common lower barrier. Their exponential error tends to zero even if the approximating masses grow, so the lower-barrier criterion identifies the limit. This addition step will remove a measure deficit in the prescribed-type existence proof.

\Needspace{15\baselineskip}
\begin{theorem}\label{thm:weighted-measure-addition}\index{Monge--Ampere measure@Monge--Amp\`ere measure!weighted addition}
Let $\theta\geq0$ be big, and suppose $\varphi\in\PSH(X,\theta)$ admits $P_{\delta\theta}(\varphi)\not\equiv-\infty$ for some $0<\delta<1$. Let $u,v\in\PSH(X,\theta)$ satisfy $u,v\sim\varphi$, and let $\lambda>0$. If
\[
 \mu=\mathrm{e}^{-\lambda u}\mathrm{MA}^{\mathrm{ext}}_\theta(u)
       +\mathrm{e}^{-\lambda v}\mathrm{MA}^{\mathrm{ext}}_\theta(v)
\]
is a finite measure, then there exists $a\in\PSH(X,\theta)$, with $a\sim\varphi$, such that
\[
 \mathrm{MA}^{\mathrm{ext}}_\theta(a)=\mathrm{e}^{\lambda a}\mu.
\]
Weights are interpreted off the pluripolar pole sets. The solution is unique in this ordinary singularity type.
\end{theorem}
\begin{proof}
Choose $S\geq\max(\sup_Xu,\sup_Xv)$. A common translation reduces to $u,v\leq0$: translating by $-S$ multiplies $\mu$ by $\mathrm{e}^{\lambda S}$ and translates the resulting solution back by $S$.

For integers $j\geq1$, set $u_j=\max(u,-j)$, $v_j=\max(v,-j)$ and
\[
 \mu_j=\mathrm{e}^{-\lambda u_j}\mathrm{MA}_\theta(u_j)
           +\mathrm{e}^{-\lambda v_j}\mathrm{MA}_\theta(v_j).
\]
For $s_j=(u_j+v_j)/2$, positivity gives
\[
 \mu_j\leq2^n \mathrm{e}^{\lambda j}\mathrm{MA}_\theta(s_j).
\]
Each measure is finite and nonzero by individual bounded noncollapse. The bounded subsolution theorem gives $a_j$ with
$\mathrm{MA}_\theta(a_j)=\mathrm{e}^{\lambda a_j}\mu_j$.
Bounded exponential order yields
\[
 a_j\leq\min(u_j,v_j)\leq0.
\]
Choose $C$ with $|u-v|\leq C$ off the common poles; then $|u_j-v_j|\leq C$. Put
\[
 D=C/2+n\lambda^{-1}\log2,\qquad b_j=s_j-D.
\]
The inequality
\[
 \mathrm{e}^{-\lambda b_j}\mathrm{MA}_\theta(b_j)
 \geq2^{-n}\mathrm{e}^{\lambda D-\lambda s_j}
            \bigl(\mathrm{MA}_\theta(u_j)+\mathrm{MA}_\theta(v_j)\bigr)
 \geq\mu_j
\]
follows from $|u_j-v_j|\leq C$. Hence $a_j\geq b_j\geq w$, where
\[
 w=(u+v)/2-D\sim\varphi
\]
is a fixed genuine potential. Extract $a_j\to a$ in $L^1(X)$. The upper barriers decrease to $\min(u,v)$, so $w\leq a\leq\min(u,v)$ and $a\sim\varphi$.

To prove convergence in capacity, write $E_j=\{\min(u,v)>-j\}$. Locality gives $\mu_j=\mu$ there. Thus
\[
 \int_{E_j}|\mathrm{e}^{a_j}-\mathrm{e}^a|\,\mathrm{MA}_\theta(a_j)
 \leq\int_X|\mathrm{e}^{a_j}-\mathrm{e}^a|\,\mathrm{d}\mu\longrightarrow0
\]
by \cref{lem:fixed-npp-integrals}. On $X\setminus E_j$ we have $a_j\leq-j$ and $a\leq-j$, whence $|\mathrm{e}^{a_j}-\mathrm{e}^a|\leq \mathrm{e}^{-j}$. Also
\[
 \mathrm{MA}_\theta(a_j)
 =\mathrm{e}^{\lambda(a_j-u_j)}\mathrm{MA}_\theta(u_j)
    +\mathrm{e}^{\lambda(a_j-v_j)}\mathrm{MA}_\theta(v_j)
 \leq\mathrm{MA}_\theta(u_j)+\mathrm{MA}_\theta(v_j).
\]
Choose a Hermitian $\Omega\geq\theta$. Write $r_n=\max(n-2,0)$. The bounded-range estimate gives
\[
 \int_{X\setminus E_j}|\mathrm{e}^{a_j}-\mathrm{e}^a|\,\mathrm{MA}_\theta(a_j)
 \leq C'\mathrm{e}^{-j}(1+j)^{r_n}\longrightarrow0.
\]
The lower-barrier criterion \cref{prop:lower-barrier-criterion} now gives $a_j\to a$ in capacity. Only polynomial control of the approximating masses was needed.

It remains to identify the equation. This can be done on bounded levels without claiming raw weak convergence prematurely. For $T>0$ choose the continuous cutoff $g_T$ from \cref{lem:weighted-extended-lsc} and put $\zeta_T=g_T(w)$. Where $\zeta_T>0$, we have $w>-2T$, hence every $a_j$ and $a$ is above $-2T$. If $j>2T$ and $x\notin E_j$, the inequalities $w(x)\leq a_j(x)\leq-j$ force $\zeta_T(x)=0$. Consequently, for all $j>2T$,
\[
 \zeta_T\mathrm{MA}_\theta(a_j)
       =\zeta_T \mathrm{e}^{\lambda a_j}\mu.
\]
Test this identity against any continuous function. On the left, replace $a_j$ by $\max(a_j,-2T)$ and use bounded capacity continuity with the fixed bounded quasi-continuous weight $\zeta_T$. On the right, fixed-measure exponential convergence gives convergence in total variation. Therefore
\[
 \zeta_T\mathrm{MA}^{\mathrm{ext}}_\theta(a)
       =\zeta_T\mathrm{e}^{\lambda a}\mu.
\]
Let $T\to\infty$. The cutoffs increase to one off the pole set of the genuine $w$; all measures ignore that set. This proves the equation and, in particular, finiteness of the limiting raw measure. Uniqueness follows from singular exponential order.
\end{proof}
